% =====================================================================
\section{The cortex: excitation and inhibition}
\fig{fig_circuit.png}{The simplest circuit of the cortex: excitatory (pyramidal) cells excite an inhibitory cell, which inhibits them in return. On the right, three molecules discussed below.}
Most cortical neurons excite one another. Such a network is unstable by itself: activity would run away to its limit. It is held in place by inhibitory cells, which receive excitation from their neighbours and inhibit them in return (fig.~4). The fastest of these contain the protein \emph{parvalbumin} and are called PV cells.
\facts
\begin{itemize}
\item About 80\,\% of cortical neurons are excitatory (\emph{pyramidal}) and about 20\,\% inhibitory.\tU\ PV cells make up roughly 40\,\% of the inhibitory cells of the mouse cortex.\tU
\item A single PV cell receives 16--34 thousand synapses; inhibition reaches a pyramidal cell with a delay of no more than 2\,ms; firing rates reach a hundred impulses per second.\tS
\item PV cells have a high energy turnover.\tS
\item Parvalbumin binds calcium and removes it quickly after an impulse.\tU\ \emph{GAD67} is the enzyme that makes GABA: less GAD67 means weaker inhibition.\tU
\item PV cells are surrounded by a shell of extracellular material (the \emph{perineuronal net}). It protects the cell from oxidants and matures during adolescence.\tS
\item PV cells set the \emph{gamma rhythm} --- oscillations of cortical activity at 30--80\,Hz. In patients with schizophrenia this rhythm is disturbed.\tU
\item The cortex operates in a regime where excitation amplifies itself and is held in check only by inhibition.\tS
\end{itemize}
\begin{ineq}
A PV cell is a fast negative feedback without which the network is unstable, and at the same time its clock generator. The price of speed is energy turnover, and therefore a flux of oxidants (section~8). This is the most heavily loaded element of the circuit.
\end{ineq}
